% ============================================================
% CAPÍTULO 1: INTRODUCCIÓN
% ============================================================

\chapter{Introducción}


% ------------------------------------------------------------
\section{Descripción del problema}
% ------------------------------------------------------------

% PÁRRAFO 1: CONTEXTO GENERAL
% Introducir los accionamientos eléctricos y su importancia.
% Hablar brevemente de la evolución desde máquinas de CC hacia
% máquinas de CA gracias al desarrollo de electrónica de potencia
% y procesadores digitales.
% Explicar que muchas aplicaciones requieren control preciso de
% velocidad y/o posición.
%
% NO entrar todavía en clasificación detallada de motores.
% NO incluir ecuaciones.


Los motores eléctricos son máquinas que transforman energia electrica en movimiento mecánico, por lo que son usadas en un sinfin de aplicaciones. Inicialmente, se utilizaban motores de corriente continua, sin embargo, con el avance de dispositivos electrónicos, controladores y procesadores digitales, se ha optado por utilizar máquinas de corriente alterna. Si bien su control es más complejo, pueden llegar a ser menos costosos, con mayor potencia, mejor respuesta dinámica,velocidad máxima, mantenimiento,entre otros factores. \cite{vierbergenPredictingImprovingBehaviour} \cite{deangeloControlParaMaquinas2004}.

A medida que ha avanzado la tecnología,la necesidad de precisión ha crecido también.
Para esto  se suele usar un sensor de posición ó velocidad para cerrar el lazo de control y maximizar la precisión de velocidad ó posición. Sin embargo, hay aplicaciones dónde por la naturaleza del sistema, se necesita un control preciso y constante sin estos sensores. 



% PÁRRAFO 2: INTRODUCCIÓN AL MOTOR PASO A PASO
% Reducir el problema general hacia los motores paso a paso.
% Explicar por qué son atractivos:
%   - Movimiento incremental.
%   - Posibilidad de controlar posición mediante excitación eléctrica.
%   - Operación tradicional a lazo abierto.
%   - Simplicidad y costo.
%   - Uso en aplicaciones de posicionamiento y baja velocidad.
%
% Solo introducir el stepper.
% La clasificación VR / PM / híbrido va en el Marco Teórico.

Para ello, salen alternativas cómo los motores paso a paso(de ahora en adelante, steppers), que pueden mejorar esta precisión a lazo abierto mediante la aplicación de pasos eléctricos.Estos motores, son motores que tienen un movimiento incremental,es decir, que tiene pequeños incrementos angulares discretos, ó pasos. Esto se realiza mediante un cambio ó excitación de las fases de este tipo de motor.
Por esto mismo, este tipo de motor se utiliza bastante, puesto que, puede comandarse en lazo abierto .Adicionalmente, la facilidad de su uso mediante drivers, librerías, y su fácil mantenimiento y costo lo hacen bastante atractivo para la industria ó proyectos de por si.



% PÁRRAFO 3: LIMITACIONES DEL CONTROL A LAZO ABIERTO
% Explicar que la posición eléctrica impuesta no garantiza que
% el rotor siga perfectamente la referencia.
%
% Introducir problemas como:
%   - Perturbaciones de carga.
%   - Variaciones paramétricas.
%   - Pérdida de pasos.
%   - Oscilaciones.
%   - Resonancias.
%   - Ripple de velocidad y/o par.
%
% Destacar que estos efectos son relevantes cuando se requiere
% movimiento suave y preciso, especialmente a baja velocidad.

Aunque estos motores tienden a presentar una alta capacidad de torque, su precisión tanto en velocidad como en posición puede variar dependiendo de las perturbaciones. Al aplicar una velocidad de pasos para un motor sin carga no es lo mismo que aplicar la misma velocidad pero con carga del 50,70 ó 100 \%.
Esto se puede traducir en una pérdida de pasos, oscilaciones, ripple en la velocidad,etc.
Una forma que se realiza para reducir este tipo de desconocimiento de parámetros ó perturbaciones, es aplicar la mayor capacidad de corriente independiente de la carga, intentando realizar un control con todo el potencial eléctrico. Esto no asegura necesariamente que no haya perdida de pasos, además de aumentar el consumo de corriente, aumentar temperatura y disminuir el tiempo de vida del motor.
Estos efectos son tema de estudio para contrarrestarse y lograr un movimiento suave y preciso, especialmente a baja velocidad, dónde las perturbaciones afectan la velocidad comaprado a alta velocidad,situación dónde se podrían hasta despreciar.


% PÁRRAFO 4: CONTROL EN LAZO CERRADO
% Introducir el control realimentado como solución a las
% limitaciones anteriores.
%
% Explicar que conocer la posición y/o velocidad real del rotor
% permite:
%   - Mejorar precisión.
%   - Rechazar perturbaciones.
%   - Evitar pérdida de sincronismo.
%   - Mejorar desempeño dinámico.
%
% Mencionar encoder/resolver únicamente como ejemplos.

Para contrarrestar esto,se suelen colocar encoders ó sensores de posición, los cuales pueden aumentar esta necesitada precisión, rechazar perturbaciones, mejorar el desempeño,etc. Pero estos sensores se traducen en mayor costo, cableado, mantenimiento, y pierde una de las ventajas de los motores paso a paso comparado a otro motor de CA, que es el poder óperar a lazo abierto.

% PÁRRAFO 5: PROBLEMA DE LOS SENSORES MECÁNICOS
% Explicar las desventajas de utilizar sensores:
%   - Mayor costo.
%   - Cableado adicional.
%   - Mayor volumen.
%   - Complejidad de integración.
%   - Mantenimiento.
%   - Confiabilidad en ciertos ambientes.
%
% Este párrafo debe conducir naturalmente hacia el concepto
% de control sensorless.


% PÁRRAFO 6: CONTROL SENSORLESS
% Introducir el control sin sensores mecánicos.
%
% Explicar conceptualmente que posición y velocidad pueden
% estimarse utilizando:
%   - Tensiones.
%   - Corrientes.
%   - Modelo matemático de la máquina.
%   - Observadores/estimadores.
%
% Se puede mencionar el EKF como una alternativa existente,
% pero SIN desarrollar todavía sus ecuaciones.

Sin embargo, hay técnicas que permiten observar y/ó estimar la velocidad de un motor de CA. Para esto, es necesario realizar una medición de algunas variables del motor, tales cómo el voltaje, corriente, algunos parámetros de la máquina y desarrollar un observador ó estimador apropiado.

% PÁRRAFO 7: PROBLEMA SENSORLESS A BAJA VELOCIDAD
% Este es uno de los puntos centrales de la problemática.
%
% Explicar:
%   - Disminución de la FEM a baja velocidad.
%   - Menor información disponible para estimar estados mecánicos.
%   - Problemas de observabilidad.
%   - Mayor sensibilidad frente a errores de modelo.
%   - Mayor efecto relativo del ruido y las perturbaciones.
%
% Respaldar estas afirmaciones con papers.
%
% Conducir hacia la idea de que el control sensorless a baja
% velocidad continúa siendo una condición desafiante.

Una de las problemáticas de observar las variables de un motor se presenta al observar con una baja velocidad mecánica, debido a la disminución de la FEM, puesto que depende de la velocidad del rotor, disminuyendo la información disponible para la estimación de los estados del motor. Por ende, las señales utilizadas presentar una menor relación señal-ruido y es más sensible a perturbaciones, ruido de mediciones y errores de modelo. Esto dificulta distinguir si las variables observadas corresponden al comportamiento real del sistema ó a ótros efectos, deteriorando la estimación de las variables. Sin embargo, hay técnicas de estimación que permiten extender este limite inferior de observabilidad, aumentando las corrientes de flujo  ó campo eléctrico, sin embargo, esto no asegura la observabilidad a velocidad zero. \cite{wangParameterRobustnessImprovement2026}.

% PÁRRAFO 8: PERTURBACIONES PERIÓDICAS Y RIPPLE
%
% IDEA:
% Hasta este punto se estableció que la estimación sensorless se vuelve
% más compleja a baja velocidad. Ahora introducir una segunda problemática:
% las perturbaciones periódicas propias de la máquina.
%
% ORDEN RECOMENDADO:
%
% 1. Explicar que el torque desarrollado por una máquina real no es
%    perfectamente constante, sino que puede contener componentes
%    periódicas.
% 2. Introducir el concepto de torque ripple como las oscilaciones
%    del torque alrededor de su valor medio.
%
% 3. Explicar que estas oscilaciones pueden tener distintos orígenes:
%       - geometría de la máquina;
%       - armónicos electromagnéticos;
%       - conmutación/excitación;
%       - interacción rotor-estator;
%       - detent torque.
%
% 4. Particularizar hacia motores paso a paso.
%    Introducir el detent torque como una perturbación periódica
%    asociada a la posición del rotor que puede existir incluso
%    cuando las fases no están energizadas.
%
% 5. Explicar la consecuencia mecánica:
%
%       perturbación periódica de torque
%                     ↓
%             oscilación de torque
%                     ↓
%             ripple de velocidad
%
% 6. Conectar esto con BAJA VELOCIDAD:
%    a bajas velocidades, estas perturbaciones pueden adquirir
%    mayor importancia relativa respecto del torque requerido
%    para mantener el movimiento.
%
% 7. Finalmente conectar con SENSORLESS:
%    si la velocidad utilizada por el controlador proviene de un
%    observador, la perturbación afecta a la planta mientras que
%    la compensación depende de variables mecánicas estimadas.
%
% IMPORTANTE:
% No desarrollar todavía la ecuación sinusoidal del detent torque.
% Eso corresponde al marco teórico/modelado.
%
% Tampoco afirmar todavía que el detent torque es "la principal"
% causa del ripple. Eso deberá justificarse bibliográficamente.
La capacidad de un motor para generar torque sin variaciones de carga es idealmente constante, pero la realidad es algo diferente. Dentro de las propiedades internas existen ciertas configuraciones propias del motor que pueden generar componentes periódicos en el torque generado por el motor. A estas oscilaciones se les llama como torque ripple, las cuales pueden ser generadas por distintas variables, cómo la geometría del motor, armónicos electromagnéticos, el tipo de conmutación, interacción rotor-stator o torque de detención. Muchas de estas perturbaciones son en función de la posición, por lo tanto,están presente independiente de si el motor está energizado ó no. Los motores stepper son capaces de generar pasos con precisión, en gran medida gracias a estas características,puesto que son construidos pensados aprovechando este tipo de diseño. 

Durante el movimiento del rotor, estas componentes periódicas de torque pueden producir oscilaciones en la velocidad del motor.Estos efectos pueden aumentar a baja velocidad, dónde las variaciones pueden ser mucho más significativas, deteriorando el comportamiendo de velocidad.Debido a la diversidad de fenómenos que pueden contribuir a estas oscilaciones,el presente trabajo se centrará específicamente en la componente
periódica asociada al \textit{cogging torque}, estudiando su efecto
sobre la velocidad del motor y su posterior compensación.

En cuestión del cogging, este se suele compensar para disminuir ó eliminar sus efectos. En la literatura, las principales formas de compensación son:


\begin{enumerate}
    \item \textbf{Métodos pasivos:} corresponden a modificaciones
    constructivas de la máquina destinadas a reducir la perturbación
    desde su origen, tales como modificaciones en la geometría del
    rotor o estator.

    \item \textbf{Métodos activos:} utilizan estrategias de control
    para reducir o compensar el efecto de la perturbación durante
    el funcionamiento del motor. Estos métodos pueden clasificarse
    a su vez en:

    \begin{enumerate}
        \item \textbf{Métodos offline:} la perturbación es caracterizada
        o identificada previamente, utilizando posteriormente esta
        información durante la operación del motor.

        \item \textbf{Métodos online:} la perturbación es estimada o
        identificada durante el funcionamiento del motor, permitiendo
        adaptar la acción de compensación a las condiciones de operación.
    \end{enumerate}
\end{enumerate}


% PÁRRAFO 9:
% Introducir de manera GENERAL las alternativas de compensación.
%
% Puedes separar conceptualmente:
%
% 1. MÉTODOS PASIVOS / CONSTRUCTIVOS
%    Modifican el diseño de la máquina para reducir la generación
%    de la perturbación.
%
%    Ejemplos:
%       - geometría del rotor/estator;
%       - modificación de dientes;
%       - skewing;
%       - optimización electromagnética.
%
% 2. MÉTODOS ACTIVOS / DE CONTROL
%    No modifican físicamente la máquina, sino que intentan
%    contrarrestar la perturbación mediante el accionamiento.
%
%    Ejemplos generales:
%       - modificación de las referencias de corriente;
%       - compensación feedforward;
%       - control resonante;
%       - observadores de perturbación.
%
% Dentro de las estrategias activas puedes introducir otra
% clasificación importante para TU tesis:
%
%       OFFLINE                     ONLINE
%          |                           |
% identificación previa       estimación/adaptación
% de la perturbación          durante el funcionamiento
%          |                           |
% LUT / mapa / perfil         observador / estimador /
% predefinido                 compensación adaptativa
%
%
% IMPORTANTE:
% En esta sección NO hacer todavía una revisión paper por paper.
%
% Solo quieres demostrar:
%
% "El problema existe y hay distintas formas de abordarlo,
% pero algunas requieren conocimiento previo de la perturbación
% o información mecánica del rotor."
%
% Esto prepara directamente el problema específico de tu tesis.

% ============================================================
\section{Antecedentes del tema}
% ============================================================

% Esta sección corresponde principalmente al ESTADO DEL ARTE.
% Aquí se utilizan los papers seleccionados para mostrar:
%
%   1. Qué se ha hecho.
%   2. Cómo se ha hecho.
%   3. Qué resultados se han obtenido.
%   4. Qué limitaciones permanecen.
%
% Debe conducir progresivamente hacia el GAP de investigación.


% ------------------------------------------------------------
\subsection{Control de motores paso a paso}
% ------------------------------------------------------------

% Revisar trabajos relacionados con:
%   - Control tradicional de steppers.
%   - Microstepping.
%   - Control de corriente.
%   - Control de velocidad.
%   - Control de posición.
%   - Control vectorial / FOC aplicado a steppers.
%   - Lazo abierto versus lazo cerrado.
%
% El objetivo NO es explicar matemáticamente cada estrategia.
% El objetivo es mostrar la evolución de las técnicas de control.


% ------------------------------------------------------------
\subsection{Control sensorless de motores paso a paso}
% ------------------------------------------------------------

% Revisar específicamente los métodos sensorless encontrados
% en la literatura.
%
% Por ejemplo:
%   - Métodos basados en FEM.
%   - Observadores de estado.
%   - EKF.
%   - Sliding Mode Observer.
%   - MRAS, si aparece en tus papers.
%   - HFI/LFI.
%   - Métodos basados en saliencia.
%
% Comparar:
%   - Qué estiman.
%   - Qué variables utilizan.
%   - Rango de velocidad.
%   - Requerimientos.
%   - Limitaciones.
%
% Esta subsección debería conducir hacia:
%
%       SENSORLESS
%           ↓
%   PROBLEMAS A BAJA VELOCIDAD


% ------------------------------------------------------------
\subsection{Control sensorless a baja velocidad}
% ------------------------------------------------------------

% Hacer zoom específicamente en la región de baja velocidad.
%
% Revisar:
%   - Problemas de observabilidad.
%   - Disminución de FEM.
%   - Saliencia.
%   - HFI.
%   - LFI.
%   - Observadores no lineales.
%   - EKF a baja velocidad.
%
% Analizar hasta qué velocidad trabajan los papers.
%
% Esto es especialmente importante para justificar por qué
% estudiar aproximadamente 1--10 rad/s (o el rango definitivo
% de la tesis) constituye un problema relevante.


% ------------------------------------------------------------
\subsection{Perturbaciones periódicas y ripple}
% ------------------------------------------------------------

% Revisar literatura relacionada con:
%   - Detent torque.
%   - Cogging torque.
%   - Torque ripple.
%   - Armónicos espaciales.
%   - Ripple de velocidad.
%
% Analizar:
%   - Origen físico.
%   - Dependencia con posición.
%   - Dependencia con velocidad.
%   - Frecuencias características.
%   - Modelos utilizados en los papers.
%
% Se pueden mencionar representaciones sinusoidales/Fourier,
% pero el desarrollo matemático detallado debe quedar para
% el Marco Teórico.


% ------------------------------------------------------------
\subsection{Métodos de reducción y compensación del ripple}
% ------------------------------------------------------------

% Revisar las estrategias encontradas para reducir ripple.
%
% Una forma útil de organizar esta sección sería:
%
% A) MÉTODOS OFFLINE
%   - Look-up tables.
%   - Mapas de compensación.
%   - Perfiles de corriente preidentificados.
%   - Identificación previa de armónicos.
%
% B) MÉTODOS ONLINE
%   - Observadores de perturbación.
%   - Estimadores de torque.
%   - Métodos adaptativos.
%   - Control resonante.
%   - Estimación armónica online.
%
% Para cada paper interesa identificar:
%   - ¿Usa sensor mecánico?
%   - ¿Es sensorless?
%   - ¿Trabaja a baja velocidad?
%   - ¿Estima la perturbación?
%   - ¿La compensación es online?
%   - ¿Reduce ripple de velocidad o de torque?
%   - ¿Requiere conocer previamente la perturbación?


% ------------------------------------------------------------
\subsection{Brecha de investigación}
% ------------------------------------------------------------

% Esta subsección debe ser el resultado de TODA la revisión anterior.
%
% La lógica esperada sería:
%
%   CONTROL DE STEPPERS
%          ↓
%   CONTROL SENSORLESS
%          ↓
%   BAJA VELOCIDAD
%          ↓
%   PERTURBACIONES PERIÓDICAS
%          ↓
%   COMPENSACIÓN DEL RIPPLE
%          ↓
%   ESTIMACIÓN/COMPENSACIÓN ONLINE
%          ↓
%        GAP
%
% El GAP debe construirse exclusivamente a partir de lo que
% realmente indiquen los papers revisados.
%
% Hipótesis preliminar del GAP:
%
% "Existe una oportunidad de desarrollar estrategias capaces
% de estimar y compensar en línea perturbaciones periódicas en
% accionamientos stepper sensorless operando a baja velocidad,
% utilizando las variables estimadas por el propio sistema de
% observación para reducir el ripple de velocidad."
%
% NO dejar esta frase definitiva hasta terminar la comparación
% de los papers seleccionados.



% ============================================================
\section{Objetivos}
% ============================================================


% ------------------------------------------------------------
\subsection{Objetivo general}
% ------------------------------------------------------------

% Debe responder directamente al GAP identificado.
%
% Versión preliminar:
%
% "Desarrollar una estrategia de control sensorless para un motor
% paso a paso que permita estimar y compensar en línea perturbaciones
% periódicas, con el propósito de reducir el ripple de velocidad
% durante la operación a baja velocidad."


% ------------------------------------------------------------
\subsection{Objetivos específicos}
% ------------------------------------------------------------

% Los objetivos específicos deberían seguir la metodología real
% de la tesis.
%
% Posible secuencia:
%
% 1. Modelar matemáticamente el motor paso a paso considerando
%    las perturbaciones periódicas relevantes.
%
% 2. Implementar una estrategia de estimación sensorless de
%    las variables mecánicas del motor.
%
% 3. Caracterizar el efecto de las perturbaciones periódicas
%    sobre la velocidad y sobre el estimador.
%
% 4. Desarrollar un método para estimar en línea la perturbación.
%
% 5. Diseñar una estrategia para compensar la perturbación estimada.
%
% 6. Integrar el estimador y el compensador dentro del sistema
%    de control sensorless.
%
% 7. Evaluar la reducción del ripple y el desempeño del sistema
%    mediante simulaciones y/o validación experimental.



% ============================================================
\section{Contribuciones de esta memoria}
% ============================================================

% Aquí NO repetir simplemente los objetivos.
%
% Los objetivos dicen:
%       "qué se pretende hacer".
%
% Las contribuciones dicen:
%       "qué aporta este trabajo respecto al estado del arte".
%
% Posibles contribuciones, sujetas a los resultados finales:
%
%   - Estrategia de estimación online de perturbaciones periódicas.
%
%   - Integración de dicha estimación con un esquema sensorless
%     basado en EKF.
%
%   - Compensación de perturbaciones sin utilizar medición mecánica
%     directa de posición o velocidad.
%
%   - Reducción del ripple de velocidad durante operación
%     a baja velocidad.
%
%   - Análisis del efecto del desfase/error del observador sobre
%     la compensación.
%
%   - Comparación cuantitativa entre operación con y sin
%     compensación.
%
% IMPORTANTE:
% Dejar únicamente las contribuciones que finalmente puedan
% demostrarse mediante resultados.



% ============================================================
\section{Organización de la memoria}
% ============================================================

% Describir brevemente qué contiene cada capítulo.
%
% Ejemplo:
%
% Capítulo 1:
% Introducción, problemática, antecedentes, brecha de investigación,
% objetivos y contribuciones.
%
% Capítulo 2:
% Fundamentos teóricos y modelo matemático del motor paso a paso.
%
% Capítulo 3:
% Desarrollo de la estrategia de control y estimación sensorless.
%
% Capítulo 4:
% Desarrollo de la estrategia de estimación y compensación
% de perturbaciones periódicas.
%
% Capítulo 5:
% Resultados, comparación de estrategias y discusión.
%
% Capítulo 6:
% Conclusiones y líneas de trabajo futuro.



% ============================================================
% CAPÍTULO 2: MARCO TEÓRICO
% ============================================================

\chapter{Marco teórico}


% ------------------------------------------------------------
\section{Motores eléctricos}
% ------------------------------------------------------------

% AQUÍ sí introducir formalmente los tipos de motores.
%
% Hacer una clasificación breve para llegar rápidamente al stepper.
% No transformar esta sección en una revisión general de todas
% las máquinas eléctricas.


% ------------------------------------------------------------
\section{Motores paso a paso}
% ------------------------------------------------------------

% Explicar:
%   - Principio de funcionamiento.
%   - Construcción.
%   - Rotor y estator.
%   - Número de polos/dientes.
%   - Paso mecánico.
%   - Excitación de fases.


% ------------------------------------------------------------
\subsection{Tipos de motores paso a paso}
% ------------------------------------------------------------

% Aquí colocar:
%   - Reluctancia variable (VR).
%   - Imán permanente (PM).
%   - Híbrido (HB).
%
% Explicar diferencias constructivas y justificar cuál se utiliza
% en esta memoria.


% ------------------------------------------------------------
\section{Modelo matemático del motor paso a paso}
% ------------------------------------------------------------

% Aquí comienza el desarrollo matemático formal.
%
% Incluir según corresponda:
%   - Modelo eléctrico.
%   - Modelo mecánico.
%   - FEM.
%   - Torque electromagnético.
%   - Inductancias.
%   - Saliencia.


% ------------------------------------------------------------
\section{Transformaciones de coordenadas}
% ------------------------------------------------------------

% Explicar:
%   - alpha-beta.
%   - d-q.
%   - Transformaciones utilizadas.
%   - Significado físico.
%   - Relación entre ángulo eléctrico y mecánico.


% ------------------------------------------------------------
\section{Control vectorial del motor paso a paso}
% ------------------------------------------------------------

% Explicar formalmente:
%   - Control de Id.
%   - Control de Iq.
%   - PI de corriente.
%   - Control de velocidad.
%   - Generación de referencias.
%   - Saturaciones.
%   - MTPA si finalmente forma parte de la estrategia.


% ------------------------------------------------------------
\section{Control sensorless}
% ------------------------------------------------------------

% Fundamentos matemáticos necesarios para comprender el estimador.
%
% Evitar repetir el estado del arte del Capítulo 1.
%
% Capítulo 1:
%   "qué métodos existen y qué hicieron otros autores".
%
% Capítulo 2:
%   "cómo funcionan matemáticamente los métodos que YO necesito".


% ------------------------------------------------------------
\subsection{Filtro de Kalman extendido}
% ------------------------------------------------------------

% Desarrollar:
%   - Modelo de estados.
%   - Predicción.
%   - Corrección.
%   - Matrices Q y R.
%   - Jacobianos.
%   - Estados estimados.
%
% Luego particularizarlo al modelo del stepper.


% ------------------------------------------------------------
\section{Perturbaciones periódicas}
% ------------------------------------------------------------

% Fundamentos físicos y matemáticos.
%
% Incluir:
%   - Cogging/detent torque.
%   - Torque ripple.
%   - Periodicidad respecto a theta_m.
%   - Frecuencia de perturbación respecto a Wm.
%   - Representación mediante series de Fourier.
%
% Aquí sí corresponde desarrollar expresiones como:
%
%       T_d = T_{dm}\sin(N_r\theta_m + \phi)
%
% y posteriormente sus armónicos si son utilizados.


% ------------------------------------------------------------
\section{Compensación de perturbaciones periódicas}
% ------------------------------------------------------------

% Fundamentos matemáticos necesarios para tu propuesta.
%
% Según la versión final:
%   - Control resonante.
%   - PR.
%   - Estimación armónica.
%   - Compensación feedforward.
%   - Compensación adaptativa.
%   - Desfase del sistema.
%
% Esta sección debería entregar las herramientas necesarias para
% entender posteriormente tu metodología propuesta.